\documentclass[11pt]{article}
\usepackage{fontspec}
\usepackage[french,english]{babel}
\usepackage{amsmath}
\usepackage{amssymb}
\usepackage{amsthm}
\usepackage{mathrsfs}
\usepackage{tikz-cd}
\usepackage{enumitem}

\newcommand{\printedpage}[1]{\marginpar{\footnotesize p.~#1}}
\newcommand{\msp}[1]{\marginpar{\footnotesize [#1]}}

% short macros for printed notation
\newcommand{\M}{\mathcal{M}}
\newcommand{\Filt}{\mathrm{Filt}}
\newcommand{\Fl}{\mathrm{Fl}}
\newcommand{\uFl}{\underline{\mathrm{Fl}}}
\newcommand{\uFilt}{\underline{\mathrm{Filt}}}
\newcommand{\card}{\operatorname{card}}
\newcommand{\Hom}{\operatorname{Hom}}
\newcommand{\Ob}{\operatorname{Ob}}
% the print has a wavy line over the expression (isomorphism classes);
% rendered here by a wide tilde
\newcommand{\wavy}[1]{\widetilde{#1}}

\begin{document}

\printedpage{14}
% head: A. Grothendieck : Dérivateurs, Ch. XV, Édition des Universités de Montpellier II et Paris VII
% foot: Version 5 novembre 2010
\msp{page 20}
first formula (3.7 bis), let us note that we have equality when $\M$ is the category of pairs $(A,B)$, $A$ a set and $B$ a subset of $A$, with the obvious arrows, and we define $\M^\pi$ (for $\pi\geq\pi_0$) as the strictly full subcategory consisting of the $(A,B)$ such that $\card A\geq\pi$.

\textsc{Proof of} 3.1. Let us first prove formula (3.8),
\[
\card\Hom(A,X)\ \leq\ s^{(r^{\pi_0})}\qquad\quad
\left\{\begin{array}{l}
\text{if } A\in\Filt^r,\ X\in\Filt^s,\\[1.5ex]
r\geq\pi_1,\ s\geq\sup(\pi_1,C_1).
\end{array}\right.
\]

(\textbf{N.B.} I shall not use $r\geq C_1$, postulated in 3.1 $1^\circ$.) I write \texttt{[``I write'' crossed out]}
\[
\begin{aligned}
A\ &=\ \varinjlim_I A_i\ ,\quad \card I\ \leq\ r^{\pi_0}\ ,\quad A_i\ \in\ \Filt^{\pi_0}\ ,\\
X\ &=\ \varinjlim_J X_j\ ,\quad \card J\ \leq\ s^{\pi_1}\ ,\quad X_j\ \in\ \Filt^{\pi_1}\ ,\qquad\qquad J\text{ large relative to }\pi_1\ .
\end{aligned}
\]
We shall have
\[
\Hom(A,X)\ \simeq\ \varprojlim_{i\in I}\ \underbrace{\Hom(A_i,\,\varinjlim_J X_j)}_{\overset{(\text{\footnotemark[6]})}{\simeq}\ \varinjlim\limits_{j\in J}\ \underbrace{\Hom(A_i,X_j)}_{\text{of cardinality }\leq C_1}}\ ,
\]
\footnotetext[6]{since the $A\in\Filt^{\pi_0}$ are $\pi_1$-accessible.}
thus for every $i\in I$
\[
\Hom(A_i,X)\ \leq\ \card(J)\times C_1\ \leq\ s^{\pi_1}\times C_1\ =\ s^{\pi_1}\ ,
\]
since $s\geq C_1$, hence $s^{\pi_1}\geq C_1$, whence
\[
\card(\,\varprojlim_{j\in I}\Hom(A_i,X))\ \leq\ (s^{\pi_1})^{\card I}\ \leq\ \underbrace{(s^{\pi_1})^{(r^{\pi_0})}}_{s^{\pi_1 r^{\pi_0}}}\ . % [sic] j\in I -> i\in I, the limit is over the index i of A_i
\]
Since $r\geq\pi_1$, we have $r^{\pi_0}\geq\pi_1$, hence $r^{\pi_0}\pi_1=r^{\pi_0}$, whence
\[
\card\Hom(A,X)\ \leq\ s^{(r^{\pi_0})}\ ,
\]
this is indeed formula (3.8). To derive from it an estimate for $\M^r/X$ (with $X\in\Filt^s$), \msp{page 21} I shall first bound above the cardinality of the isomorphism classes of objects of $\M^r/X$, \textit{i.e.} of the arrows up to isomorphism (in $\M^r$)
\[
A\ \longrightarrow\ X\ ,\qquad\quad X\text{ fixed in }\Filt^s,\ A\text{ variable in }\M^r\ .
\]
We shall have
\[
A\ =\ \varinjlim_{i\in I} A_i\qquad\quad\text{with } A_i\in\Filt^{\pi_0},\ I\text{ ordered set of cardinality}\leq r^{\pi_0}\ .
\]

\printedpage{15}
% head: A. Grothendieck : Dérivateurs, Ch. XV, Édition des Universités de Montpellier II et Paris VII
% foot: Version 5 novembre 2010
For fixed $I$, the cardinality of the set of direct systems indexed by $I$, in $\M^{\pi_0}/X$, up to isomorphism, can be bounded above as follows: for every $i\in I$, one must choose $A_i$ in $\M^{\pi_0}$ (but only its class in $\wavy{\M^{\pi_0}}$ matters), so this is bounded above by card $\wavy{\M^{\pi_0}}\ \leq C_0$~\textsuperscript{(}\footnote[7]{Here we have put $C_0=\card\wavy{\uFl^{\pi_0}}$ , hence $C_0\leq C_1$.}\textsuperscript{)}, and for all the $i$, by $\card I\times C_0\leq r^{\pi_0}\times C_0$, and if $r\geq C_0$, hence $r^{\pi_0}\geq C_0$, we find an upper bound on the possible choices of the system of \textit{objects} $A_i$ in $\M^\pi$, up to isomorphism, given by $r^{\pi_0}$. If in addition we specify the transition morphisms $A_i\longrightarrow A_j$ for $i<j$, this gives for each comparable pair, an element of $I\times I$, a choice bounded above by $C_0$, hence again an upper bound of $\ \underbrace{\card(I\times I)}_{=(\card I)^2=\card I}\cdot C_0\leq r^{\pi_0}\cdot C_0=r^{\pi_0}$. Thus the set of direct systems in question, up to isomorphism, for fixed $I$, has cardinality $\leq r^{\pi_0}\cdot r^{\pi_0}=r^{\pi_0}$. Finally, the \textit{homomorphisms} of a given direct system into $X$ are bounded above by $\card I\times s^{(r^{\pi_0})}$ (by (3.8), already proved), hence by
\[
r^{\pi_0}\cdot s^{(r^{\pi_0})}\ =\ s^{(r^{\pi_0})}\qquad\quad(\text{since } r^{\pi_0}\leq s^{(r^{\pi_0})})\ .
\]

Altogether, for fixed $I$, the set of direct systems up to isomorphism, indexed by $I$, with values in $\M^r/X$, is \msp{page 22} of cardinality bounded above by $\quad\underbrace{r^{\pi_0}}_{\text{\scriptsize\begin{tabular}{@{}c@{}}number of\\ direct systems\\ in $\M^\pi$, up to\\ isomorphism\end{tabular}}}\quad\cdot\quad\underbrace{s^{(r^{\pi_0})}}_{\text{\scriptsize\begin{tabular}{@{}c@{}}number of\\ arrows from\\ such a direct system\\ into $X$\end{tabular}}}\quad=s^{(r^{\pi_0})}$.

But we must now let $I$ vary, up to isomorphism, in the category of ordered sets (filtering \textit{etc.}) of cardinality $\leq\lambda=r^{\pi_0}$. We see at once that this set of isomorphism classes (which can be interpreted as a quotient set of a certain set of subsets of a fixed set $I_0$ of cardinality $\lambda$) has cardinality $\leq 2^\lambda$. We thus finally find
\[
\card\wavy{\M^r/X}\ \ \leq\ s^{(r^{\pi_0})}\cdot 2^{(r^{\pi_0})}\ =\ s^{(r^{\pi_0})}\ ,
\]
on condition that
\[
r\ \geq\ \sup(\pi_1,C_0)\ ,\qquad\quad s\ \geq\ \sup(\pi_1,C_1)\ .
\]
We must prove the similar inequality for card $\wavy{\uFl^r\,\M}$ \ (3.9). Now, since for $A$, $B$ in $\M^r$, we have, if $r\geq\sup(\pi_1,C_1)$,
\[
\card\Hom_{\M}(A,B)\ \leq\ r^{(r^{\pi_0})}
\]
by virtue of (3.8) where we put $s=r$, we find that
\[
\begin{aligned}
\card\wavy{\uFl(\M^r/X)}\ &\leq\ \Bigl(\card\wavy{\M^r/X}\,\Bigr)^2\cdot r^{(r^{\pi_0})}\\
&\leq\ s^{(r^{\pi_0})}\cdot r^{(r^{\pi_0})}\\
&=\ s^{(r^{\pi_0})}\qquad\text{if } s\geq r\ .
\end{aligned}
\]
(We shall look below at the cases $s\leq r$~\textsuperscript{(}\footnote[8]{\textit{cf.} cor. 2, page 23.}\textsuperscript{)}). This is indeed formula (3.9).

\printedpage{16}
% head: A. Grothendieck : Dérivateurs, Ch. XV, Édition des Universités de Montpellier II et Paris VII
% foot: Version 5 novembre 2010
\msp{page 23}
Let us turn to the category $\M^r$ itself. The argument on page 21 shows that the set of isomorphism classes of objects of $\M^r$ indexed by a fixed $I$, of cardinality $\leq r^{\pi_0}$, as $A=\varinjlim\limits_i A_i$ with $A_i\in\Filt^{\pi_0}$, has cardinality bounded above by $r^{\pi_0}$ (if $r\geq C_0$). If we let $I$ vary, we find an upper bound of $r^{\pi_0}\cdot 2^{(r^{\pi_0})}=2^{(r^{\pi_0})}$. We conclude from this as before that the set $\ \wavy{\Fl\,\M^r}$ has cardinality bounded above by
\[
(2^{(r^{\pi_0})})^2\times r^{(r^{\pi_0})}\qquad\quad(\text{if } r\geq\sup(\pi_1,C_1))\ ,
\]
hence by $r^{(r^{\pi_0})}$. I say that we have
\[
r^{(r^{\pi_0})}\ =\ 2^{(r^{\pi_0})}\ ,
\]
indeed, we have $\geq$, and to prove $\leq$, we note that $r<2^r$, hence $r^{(r^{\pi_0})}\leq(2^r)^{(r^{\pi_0})}=2^{r\cdot r^{\pi_0}}=2^{(r^{\pi_0})}$, O.K.! This thus establishes 3.7. This argument in fact proves the

\textbf{Lemma 3.2.} \textit{Let $r$ be a cardinal, $r'$ a cardinal $\geq r$. Then we have}
\[
r^{r'}\ =\ 2^{r'}\qquad\quad(=s^{r'}\ \textit{for all}\ 2\leq s\leq r)\ .
\]

\textbf{Corollary 1.} \textit{Let $r\geq\sup(\pi_1,C_1)$, and $A,X\in\Filt^\pi$. Then}
\[
\card\Hom_{\M}(A,X)\ \leq\ 2^{(r^{\pi_0})}\qquad\quad(=2^r\ \textit{if}\ r=2^c,\ \textit{with}\ c\geq\pi_0)\ . \tag*{(3.10)}
\]

We apply (3.8), which gives the upper bound $r^{(r^{\pi_0})}$, and lemma 3.2.

We shall note that we find the same upper bound for the cardinality of $\wavy{\uFl(\M^r)}$ as a whole (3.7)! Likewise, (3.9) gives:

\textbf{Corollary 2.} \textit{Let $r\geq\sup(\pi_1,C_1)$, and $X\in\Ob\M^r$. Then}
\[
\card\wavy{\uFl(\M^r/X)}\ \leq\ 2^{(r^{\pi_0})}\qquad\quad(=2^r\ \textit{if}\ r=2^c,\ \textit{with}\ c\geq\pi_0)\ . \tag*{(3.11)}
\]
so this is formula (3.9) for $s\leq r$, left pending.

\msp{page 24}
Let us return to the cardinal bound for the $X_\alpha(f)$, where $f : X\longrightarrow Y$ is in $\Fl^s(\M)$, \textit{i.e.} $X,Y\in\Filt^s$. We assume moreover $C_0\subset\Fl^r(\M)$,
\[
\begin{aligned}
&C_0\ \subset\ \Fl^r(\M)\ =\ \Fl(\M^r)\ ,\qquad\quad(f:X\longrightarrow Y)\in\Fl^s(\M)\ ,\ \ \textit{i.e.}\ X,Y\in\Filt^s\ ,\\
&\text{with } s\geq r\ .
\end{aligned} \tag*{(3.10)}
\]

\texttt{[Twice (3.10).]} There is a stupid difficulty, which is that $\M^r$ is equivalent to a small category, but is not itself small. Even if $C_0$ contained only a single isomorphism class of arrows of $\M^r$, its cardinality could be (for fixed $r$) arbitrarily large. Thus, in what follows, I must assume the condition of ``cardinal non-plethoricity'' for $C_0$:
\[
\parbox{0.72\linewidth}{$\forall\,i\in C_0$, denoting by $C_0(i)$ the set of $i'\in C_0$ such that $i\simeq i'$ in $\Fl(\M)$, I require that $\card C_0(i)\leq 2^{(r^{\pi_0})}$ (which is also the value of the upper bound for the cardinality of $\wavy{\uFl^r(\M)}=\wavy{\uFl(\M^r)}$\,).} \tag*{(3.11)}
\]

\printedpage{17}
% head: A. Grothendieck : Dérivateurs, Ch. XV, Édition des Universités de Montpellier II et Paris VII
% foot: Version 5 novembre 2010
\texttt{[Twice (3.11).]} This thus implies that we have $\card C_0\leq(2^{(r^{\pi_0})})^2=2^{(r^{\pi_0})}$,
\[
\card C_0\ \leq\ 2^{(r^{\pi_0})}\ , \tag*{(3.12)}
\]
which is equivalent to (3.11) (given $C_0\subset\Fl^r(\M)$, postulated in (3.10)).

It follows by (3.8) that we have
\[
\card\Delta(C_0,f)\ \leq\ \underbrace{\card(C_0)}_{\leq 2^{(r^{\pi_0})}}\cdot s^{(r^{\pi_0})}\cdot s^{(r^{\pi_0})}\ \leq\ s^{(r^{\pi_0})}\ .
\]

On the other hand, the $X_D$ whose amalgamated sum under $X$ we take, to obtain $X_1(f)$, are in $\M^s$, if we make the following assumption. \msp{page 25}
\[
\left|\ \parbox{0.72\linewidth}{The $\M^\pi=\uFilt^\pi$ are stable under simple amalgamated sums, and under filtering $\varinjlim$ of cardinality $\leq\pi$ (not necessarily large relative to $\pi_0$) ,}\right. \tag*{(3.13)}
\]
taking into account that $r\leq s$, hence $\M^r\subset\M^s$, hence $A$, $B$, $X$ in $X_D=B\amalg_A X$ are in $\M^s$. The \textit{finite} partial amalgamated sums of the $X_D$ are therefore also in $\M^s$, by induction on the cardinality of the index set. The infinite amalgamated sum thus appears as a filtering $\varinjlim$ of objects of $\M^s$, the index set being the set $\mathfrak{P}^*_{\mathfrak{f}}\Delta(C_0,f)$ of nonempty finite subsets of $\Delta(C_0,f)$. This set has cardinality bounded above by $s'$, if $s'$ is an \textit{infinite} cardinal such that $\card\Delta(C_0,f)\leq s'$. We can therefore take $s'=s^{(r^{\pi_0})}$. We have therefore proved, taking into account (3.13) applied to $\M^{s'}$, the % ?? uncertain: symbol in X_D = B \amalg_A X -- \amalg or a doubled \sqcup; subscript of \mathfrak{P}^* -- fraktur f or roman f

\textbf{Lemma 3.3.} \textit{We assume} \textup{(3.10), i.e.} $C_0\subset\Fl^r(\M)$, $f\in\Fl^s(\M)$, \textit{with} $s\geq r\geq\sup(\pi_1,C_1)$, \textit{and moreover that $C_0$ satisfies the non-plethoricity condition} \textup{(3.11)} \textup{(}\textit{which we can always assume, at the cost of reducing $C_0$ without changing $\widetilde{C_0}=\Phi$ and $C_{0*}=\Psi$}\textup{)}\textit{, and that the cardinal filtration satisfies the stability condition (harmless!)} \textup{(3.13)}\textit{. Put}
\[
s'\ =\ s^{(r^{\pi_0})}\ . \tag*{(3.14)}
\]
\textit{Then $X_1(C_0,f)$ is in $\M^{s'}$. More generally and more prettily, for fixed $C_0$ and $r$} \textup{(}\textit{with $C_0\subset\Fl^r(\M)$ and satisfying the non-plethoricity condition} \textup{(3.11)} \textit{relative to $r$}\textup{)}\textit{, for every cardinal} \msp{page 26} \textit{of the form}
\[
s'\ =\ s^t\ ,\qquad\quad\textit{with}\ t\geq r^{\pi_0},\ s\geq 2 \tag*{(3.15)}
\]
\textit{(whence $s'\geq 2^{(r^{\pi_0})}>r^{\pi_0}\geq r$), $\uFl^{s'}(\M)$ is stable under the functor $p_1^{C_0}:\uFl(\M)\longrightarrow\uFl(\M)$.}

Indeed, if we apply the first part of the lemma to the pair $(r,s')$, we find that $X_1(f)$ is in $\M^{s''}$, with $s''=s'^{(r^{\pi_0})}$. But
\[
s''\ =\ s'^{(r^{\pi_0})}\ =\ (s^t)^{(r^{\pi_0})}\ =\ s^{tr^{\pi_0}}\ \underset{\substack{\vert\\ \text{\scriptsize since $t\geq r^{\pi_0}$,}\\ \text{\scriptsize hence $tr^{\pi_0}=t$}}}{=}\ s^t\ =\ s'\ ,
\]
\printedpage{18}
% head: A. Grothendieck : Dérivateurs, Ch. XV, Édition des Universités de Montpellier II et Paris VII
% foot: Version 5 novembre 2010
so we see that $f\in\Fl^{s'}$ implies $X_1(f)\in\M^{s'}$, or what amounts to the same thing, $p_1(f)\in\Fl^{s'}$.

\textbf{Corollary.} \textit{Let $C_0$, $r$, $s'=s^t$ $(t\geq r^{\pi_0},\ s\geq 2)$ be as above, and let $\alpha$ be an ordinal number such that}
\[
\card\alpha\ \leq\ s'\ =\ s^t\ .
\]
\textit{Then $\uFl^{s'}(\M)$ is stable under the functor $p_\alpha=p_\alpha^{C_0}$.}

We see this at once by transfinite induction, taking into account lemma 3.3, and the stability assumption (3.13).

To summarize:

\textbf{Theorem 2.} \textit{Let $\M$ be a category stable under small filtering $\varinjlim$, and endowed with a cardinal filtration} $(\Filt^\pi)_{\pi\geq\pi_0}$\textit{, such that the} $\Filt^\pi$ \textit{are stable under filtering $\varinjlim$ of cardinality $\leq\pi$ (not necessarily large relative to $\pi_0$), and under the simple amalgamated sums that exist in $\M$. We assume that the elements of} $\Filt^{\pi_0}$ \textit{are accessible, and we denote by $\pi_1$ the smallest cardinal $\geq\pi_0$, such that the objects of} $\Filt^{\pi_0}$ \textit{are $\pi_1$-accessible. We put} \msp{page 27}
\[
c_1\ =\ \card\Fl(\M^{\pi_1})\ =\ \card\Fl^{\pi_1}(\M)
\]
\textit{(where for every $\pi\geq\pi_0$, $\M^\pi$ denotes the strictly full subcategory of $\M$ defined by} $\Filt^\pi$\textit{, and} $\Fl^\pi(\M)=\Fl(\M^\pi)$\textit{).}

\textit{Let $C_0$ be a small subset of} $\Fl(\M)$\textit{, consisting of cosquarable arrows, and put}
\[
\Phi\ =\ \widetilde{C_0}\ =\ \Psi^*\ ,\qquad\quad\textit{where}\ \ \Psi=C_{0*}=\Phi_*\ .
\]
\textit{Let $r$ be a cardinal} $\geq\sup(\pi_1,c_1)$\textit{, and such that we have $C_0\subset\Fl^r(\M)$. We assume that $C_0$ is irredundant (}i.e.\textit{ two isomorphic objects of $C_0$ are equal), or more generally, that $C_0$ satisfies the very harmless non-plethoricity condition} \textup{(3.11)} \textit{relative to $r$. We consider the transfinite sequence of functors $p_\alpha=p_\alpha^{C_0}$ on $\uFl(\M)$} \textup{(}\textit{cf.} \textup{th. 1)}\textit{, giving rise to the functorial factorizations of an $f\in\uFl(\M)$, $f:X\longrightarrow Y$, as}
\[
X\ \xrightarrow{i_\alpha(f)}\ X_\alpha(f)\ \xrightarrow{p_\alpha(f)}\ Y\ . \tag*{(3.16)}
\]

\textit{We then have the following:}

$1^\circ$) \textit{We have (for the record) $i_\alpha(f)\in\Phi$ for every $\alpha$, and moreover,} if $\alpha$ is large relative to $r^{\pi_0}$, \textit{we also have $p_\alpha(f)\in\Psi$. If $r'$ is a cardinal $>r^{\pi_0}$ (e.g. the successor cardinal of $r^{\pi_0}$), we can take for $\alpha$ the first ordinal whose cardinality is equal to $r'$.}

$2^\circ$) \textit{For every cardinal $s$ of the form}
\[
s\ =\ c^t\qquad\quad\textit{with}\ \ t\geq r^{\pi_0},\ c\geq 2
\]
\textit{(thus $s=2^t$ if we take $c\leq t$,} cf. \textit{lemma 3.2), $\uFl^s(\M)$ is stable under the $p_\alpha$ such that} $\card\alpha\leq s$\textit{. If we take e.g. for $\alpha$ the first ordinal whose cardinality is $>r^{\pi_0}$, then $\uFl^s(\M)$ is stable under $p_\alpha$, which by $1^\circ$) thus realizes a functorial factorization} \textup{(3.16)} \textit{with $i_\alpha(f)\in\Phi$, $p_\alpha(f)\in\Psi$.}

\end{document}
